\documentclass[11pt]{article}
\usepackage[a4paper,margin=1in]{geometry}
\usepackage{amsmath,amssymb,amsthm,mathtools}
\usepackage{microtype}
\usepackage{enumitem,xurl}
\usepackage{xcolor}
\usepackage{authblk}
\usepackage{hyperref}
\usepackage[nameinlink,capitalise,noabbrev]{cleveref}
\renewcommand{\Affilfont}{\small}
\renewcommand{\Authfont}{\normalsize}
\hypersetup{
  colorlinks=true,
  linkcolor=blue!55!black,
  citecolor=blue!55!black,
  urlcolor=blue!55!black
}
\numberwithin{equation}{section}
\newtheorem{theorem}{Theorem}[section]
\newtheorem{lemma}[theorem]{Lemma}
\newtheorem{proposition}[theorem]{Proposition}
\newtheorem{corollary}[theorem]{Corollary}
\newtheorem{remark}[theorem]{Remark}
\usepackage{etoolbox}
\AtBeginEnvironment{lemma}{\crefalias{theorem}{lemma}}
\AtBeginEnvironment{proposition}{\crefalias{theorem}{proposition}}
\AtBeginEnvironment{corollary}{\crefalias{theorem}{corollary}}
\AtBeginEnvironment{remark}{\crefalias{theorem}{remark}}
\newcommand{\Tr}{\operatorname{Tr}}
\newcommand{\one}{\mathbf 1}
\newcommand{\R}{\mathbb R}
\newcommand{\norm}[1]{\left\lVert #1\right\rVert}
\newcommand{\ip}[2]{\left\langle #1,#2\right\rangle}
\newcommand{\pos}[1]{\left(#1\right)_+}
\newcommand{\dd}{\,\mathrm d}
\title{Goodman-type lower bounds for odd cycles II: The intermediate regime}
\author[1,2]{Taeyoung~Kim}
\author[3]{Jineon~Baek}
\author[4]{Seonghyuk~Im}
\author[2]{Joonkyung~Lee}
\author[1,2]{Hongseok~Yang}
\affil[1]{School of Computing, KAIST, Korea}
\affil[2]{School of Computational Sciences, Korea Institute for Advanced Study (KIAS), Korea}
\affil[3]{June E Huh Center for Mathematical Challenges, Korea Institute for Advanced Study (KIAS), Korea}
\affil[4]{Center for AI and Natural Science, Korea Institute for Advanced Study (KIAS), Korea}
\affil[ ]{\newline\texttt{taeyoungkim21@kaist.ac.kr}, \texttt{{jineon.kias}@gmail.com}, \texttt{seonghyuk@kias.re.kr}, \texttt{joonkyunglee@kias.re.kr}, \texttt{hongseokyang@kias.re.kr}}
\date{}
\begin{document}
\maketitle
\begin{abstract}
Goodman's triangle bound suggests the polynomial lower bound
\[
 t(C_m,W)\ge p^m-p(1-p)^{m-1}
\]
for odd cycles. We prove this inequality for every graphon $W$ of edge
density $0\le p\le2/3$ and every odd integer $m\ge3$. Together with
the dense case in the companion paper, this establishes the bound at
every edge density.

The nontrivial range is $1/2<p\le2/3$. We use an excursion expansion
to separate the constant contribution from an odd spectral moment.
A Hilbert--Schmidt bound allows at most one eigenvalue to exceed the
complementary density $q=1-p$. Comparing its eigenfunction with its
absolute value forces a quantitative amount of degree variation,
whose excursion contributions compensate for the excess spectral
moment. A scalar inequality for truncated squares, combined with
Taylor's integral formula, proves this comparison simultaneously for
all odd lengths at least five.
\end{abstract}

\section{Introduction}\label{sec:intro}\label{sec:reductions}
The minimum density of an odd cycle at a prescribed edge density is
a natural instance of the subgraph-density problem. For triangles,
Goodman's bound gives the polynomial $2p^2-p$. The companion
paper~\cite{kim2026goodman} proves that the corresponding polynomial
$p^m-p(1-p)^{m-1}$ is a lower bound for every odd cycle $C_m$ when
$p\ge2/3$. At each density $p=1-1/k$, with an integer $k\ge3$,
the balanced complete $k$-partite graphon attains this bound.
The present paper treats the remaining edge densities.

We use the graphon formulation of subgraph densities
\cite{lovasz2006limits,lovasz2012large}. A \emph{graphon} is a
symmetric measurable function $W\colon[0,1]^2\to[0,1]$.
For a finite simple graph $F$, its \emph{homomorphism density} in
$W$ and the \emph{edge density} of $W$ are
\[
 t(F,W)=\int_{[0,1]^{V(F)}}\prod_{ij\in E(F)}W(x_i,x_j)
       \prod_{i\in V(F)}\dd x_i,
 \qquad p=\int_{[0,1]^2}W(x,y)\dd x\dd y.
\]
Our main result is the following.
\begin{theorem}[Odd-cycle Goodman bound for $p\le2/3$]\label{thm:main}
Let $m\ge3$ be odd, and let $W$ be a graphon of edge density
$0\le p\le2/3$. Then,
\begin{equation}\label{eq:target}
 t(C_m,W)\ge p^m-p(1-p)^{m-1}.
\end{equation}
\end{theorem}
The endpoint $p=2/3$ is also covered by
\cite[Theorem~1.1]{kim2026goodman}. Thus the two results together
give the odd-cycle Goodman bound for every $0\le p\le1$.
The proof below is self-contained, including the excursion expansion
shared with the dense argument.

\paragraph{Proof strategy.}
The dense proof adds the cycle densities of $W$ and its complement
so that their odd spectral moments cancel. It then bounds the
remaining excursion terms by a complementary-path correction.
Here we retain the odd spectral moment of the complement. The
density restriction implies that at most one eigenvalue can exceed
the level allowed by the desired bound. Such an eigenvalue forces
nonconstant degrees, and we compare the resulting excursion terms
with the excess in the spectral moment.

The proof has two distinct parts. The graphon argument isolates the
exceptional eigenfunction and derives a constraint on its coupling
to the constant direction. The scalar argument proves that this
constraint supplies a sufficient correction for every odd cycle
length. Its dependence on the length is removed by expressing powers
as positive integrals of truncated squares. The full reduction is
explained after the block decomposition in \Cref{sec:blocks}.

\paragraph{Organisation.}
We first handle the elementary cases and reduce to step graphons.
\Cref{sec:blocks,sec:excursions} establish the block decomposition
and excursion expansion. \Cref{sec:isolate} isolates the exceptional
eigenvalue, and \Cref{sec:compensation} derives the coupling constraint.
\Cref{sec:integration,sec:scalar} reduce the moment comparison to a
scalar inequality and prove it. \Cref{sec:completion} combines these
results to prove \Cref{thm:main}.

\subsection{Elementary cases and step approximation}

Put $q=1-p$. For a real number $r$, write $r_+=\max\{r,0\}$.
\begin{lemma}[Edge densities at most $1/2$]\label{lem:low-density}
If $p\le1/2$, then the bound \eqref{eq:target} holds for every odd $m\ge3$.
\end{lemma}
\begin{proof}
Since $p\le q$, we have
$p^{m}-pq^{m-1}=p(p^{m-1}-q^{m-1})\le0$, whereas $t(C_m,W)\ge0$.
\end{proof}

\begin{lemma}[Goodman's triangle bound]\label{lem:triangle}
For every graphon $W$ of edge density $p$,
$t(C_3,W)\ge p^3-pq^2$.
\end{lemma}
\begin{proof}
Let $d_W(x)=\int_0^1 W(x,z)\dd z$. For $a,b\in[0,1]$,
$(1-a)(1-b)\ge0$ gives $ab\ge a+b-1$. Applying this to
$W(x,z)$ and $W(y,z)$ and integrating in $z$ gives
\[
 \int W(x,z)W(y,z)\,dz\ge d_W(x)+d_W(y)-1.
\]
Multiplying by $W(x,y)\ge0$ and integrating in $x,y$ yields
\begin{align*}
 t(C_3,W)
 &\ge\int_{[0,1]^2}W(x,y)(d_W(x)+d_W(y)-1)\dd x\dd y\\
 &=2\int_0^1d_W(x)^2\dd x-p
 \ge2p^2-p=p^3-pq^2.
\end{align*}
Here symmetry gives the two equal squared-degree integrals, and
$\int_0^1(d_W-p)^2\dd x\ge0$ gives
$\int_0^1d_W^2\dd x\ge p^2$.
\end{proof}
This is the standard proof of Goodman's bound~\cite{goodman1959sets}.

\begin{lemma}[Reduction to step graphons]\label{lem:step-reduction}
Every graphon $W$ of edge density $p$ admits a sequence of step
graphons $W_N$ of the same density such that $W_N\to W$ in $L^1$
and $t(C_m,W_N)\to t(C_m,W)$ for every $m\ge3$.
Consequently, it suffices to prove \Cref{thm:main} for step graphons.
\end{lemma}
\begin{proof}
For an integer $N\ge1$, partition $[0,1]$ into $2^N$ intervals
$I_1,\ldots,I_{2^N}$ of length $2^{-N}$, and define
\[
 W_N(x,y)=2^{2N}\int_{I_i\times I_j}W(u,v)\dd u\dd v
 \qquad((x,y)\in I_i\times I_j).
\]
The averages lie in $[0,1]$, and symmetry of $W$ gives symmetry
of $W_N$. Moreover,
\[
 \int_{[0,1]^2}W_N
  =\sum_{i,j=1}^{2^N}\int_{I_i\times I_j}W=p.
\]
To prove $L^1$ convergence, fix $\varepsilon>0$ and choose a
function $S$ constant on the rectangles of a dyadic partition
with $\norm{W-S}_1<\varepsilon$. Such functions are dense in
$L^1([0,1]^2)$. For all sufficiently large $N$, averaging on the
level-$N$ partition fixes $S$. On each averaging cell the absolute
value of the average of $W-S$ is at most the average of $|W-S|$.
Summing over the cells therefore gives
\[
 \norm{W_N-W}_1\le\norm{W_N-S}_1+\norm{S-W}_1
  \le2\norm{W-S}_1<2\varepsilon.
\]
Finally, subtract the two products defining the cycle densities by
replacing one edge factor at a time. Each resulting term contains
one factor $W_N-W$, while every other factor lies in $[0,1]$.
Integrating its absolute value bounds the term by $\norm{W_N-W}_1$.
There are $m$ terms, so
\begin{equation}\label{eq:continuity}
 |t(C_m,W_N)-t(C_m,W)|\le m\norm{W_N-W}_1\longrightarrow0.
\end{equation}
The right-hand side of the desired bound depends only on the common
edge density, so the bound passes to the limit.
\end{proof}
This averaging argument also appears in
\cite[Lemma~2.1 and Appendix~A]{kim2026goodman}.
From now on assume
\begin{equation}\label{eq:range}
 \frac12<p\le\frac23,\qquad \frac13\le q<\frac12,
\end{equation}
and fix a step graphon $W$. Its partition has finitely many cells
$B_1,\ldots,B_N$, each of positive measure. We allow unequal cell
measures, although equal intervals suffice for the approximation.
Write $E$ for the finite-dimensional space of real functions that
are constant on these cells.

\section{Matrix reduction and spectral bounds}\label{sec:blocks}
We first express the graphon quantities as finite-dimensional
operator quantities. This identifies the odd spectral moment that
must be bounded and the degree fluctuation that supplies the
excursion correction.

A \emph{step kernel} on the partition is a symmetric real-valued
function constant on each $B_i\times B_j$. For such a kernel $K$,
define
\[
 (T_Kf)(x)=\int_0^1K(x,y)f(y)\dd y,
 \qquad \ip fh=\int_0^1 f(x)h(x)\dd x.
\]
The image of $T_K$ lies in $E$. If $f\perp E$, its integral on
each cell is zero, so $T_Kf=0$. We therefore regard $T_K$ as an
operator on $E$ and take all traces on this finite-dimensional
space. The pointwise functions $|f|$, $f_+=\max\{f,0\}$, and
$f_-=\max\{-f,0\}$ also belong to $E$ whenever $f\in E$.

Put $U=1-W$, and write $\one$ for the constant function equal to
one. It has norm one. Let $E_0=E\cap\one^\perp$, the subspace
of functions with integral zero. Define the orthogonal projections
onto the constants and onto $E_0$ by
\[
 \Pi_{\mathrm c}f=\ip\one f\one,
 \qquad \Pi_0f=f-\ip\one f\one.
\]
Relative to $E=\R\one\oplus E_0$, define
\[
 g=T_U\one-q\one,\qquad A=\Pi_0T_U|_{E_0}.
\]
For $f,h\in E_0$, this definition gives
$\ip f{Ah}=\ip f{T_Uh}$. Thus $A$ is the operator representing
the quadratic form of $T_U$ on the mean-zero subspace; the
projection is needed because $T_U$ need not preserve that subspace.
We also write $d_W(x)=\int_0^1W(x,y)\dd y$ for the degree function.
\begin{lemma}[Block decomposition and degree variance]\label{lem:blocks}
The operators have the block forms
\begin{equation}\label{eq:blocks}
 T_U=\begin{pmatrix}q&g^*\\g&A\end{pmatrix},\qquad
 T_W=\begin{pmatrix}p&-g^*\\-g&-A\end{pmatrix}.
\end{equation}
Moreover $A$ is self-adjoint, $g=p-d_W$, and
\[
 \norm g^2=\int_0^1(d_W(x)-p)^2\,dx.
\]
In particular, $g=0$ exactly when $W$ has constant degree.
\end{lemma}
\begin{proof}
The edge density of $U$ is $q$, so
\[
 \ip\one g=\ip\one{T_U\one}-q
  =\int_{[0,1]^2}U-q=0.
\]
Hence $g\in E_0$. The first column of $T_U$ is given by
$T_U\one=q\one+g$. For $f\in E_0$, symmetry of $U$ gives
$\ip\one{T_Uf}=\ip{T_U\one}f=\ip gf$, which determines the
first row. Here $g^*$ denotes the linear functional $f\mapsto\ip gf$.
The remaining block is $A$ by definition, and
$\ip f{Ah}=\ip{Af}h$ proves its self-adjointness.
Moreover,
\[
 (T_W+T_U)f=\int_0^1 f(y)\dd y\,\one
  =\Pi_{\mathrm c}f.
\]
Subtracting the matrix of $T_U$ from that of $\Pi_{\mathrm c}$
gives the claimed matrix of $T_W$. Finally,
\[
 g(x)=\int_0^1(1-W(x,y))\dd y-q=p-d_W(x).
\]
Squaring and integrating proves the variance identity. Its right-hand
side vanishes exactly when $d_W=p$ almost everywhere.
\end{proof}
This is the block form of \cite[Eq.~(2.2)]{kim2026goodman}, there
written for cells of equal measure. The eigenvalues of $A$ describe
walks that stay in the mean-zero subspace, while $g$ connects that
subspace to the constants.

\begin{lemma}[Trace formula]\label{lem:tracecycle}
For a step kernel $K$ and an integer $m\ge2$,
\begin{equation}\label{eq:tracecycle}
 \Tr(T_K^m)=\int_{[0,1]^m}K(x_1,x_2)K(x_2,x_3)\cdots K(x_m,x_1)\,dx_1\cdots dx_m .
\end{equation}
For $m\ge3$ the right side is $t(C_m,K)$; for $m=2$ it is $\int K^2$.
\end{lemma}
\begin{proof}
If a cell $B_i$ has measure $b_i$, the matrix of $T_K$ in the
orthonormal basis $\one_{B_i}/\sqrt{b_i}$ has entries
$K_{ij}\sqrt{b_ib_j}$. Thus
\[
 \Tr(T_K^m)=\sum_{i_1,\ldots,i_m}
 K_{i_1i_2}\cdots K_{i_mi_1}b_{i_1}\cdots b_{i_m}.
\]
Each summand is the integral over
$B_{i_1}\times\cdots\times B_{i_m}$ in \eqref{eq:tracecycle}.
\end{proof}
The trace interpretation of cycle densities is standard; see
\cite{lovasz2012large}.

\begin{lemma}[The regular case]\label{lem:regular-case}
If $g=0$, then $t(C_m,W)=p^m-\Tr(A^m)$ for every odd $m\ge3$.
\end{lemma}
\begin{proof}
The block matrix of $T_W$ is $p\oplus(-A)$. By the trace formula,
$t(C_m,W)=p^m+\Tr((-A)^m)=p^m-\Tr(A^m)$.
\end{proof}
For a nonregular graphon, excursions through the off-diagonal blocks
supply additional terms. Their sign, and their relation to a large
positive eigenvalue of $A$, are the two issues in the proof.

\begin{lemma}[Hilbert--Schmidt bound]\label{lem:HS}
One has
\begin{equation}\label{eq:HS}
 \Tr(A^2)+2\norm g^2\le pq.
\end{equation}
In particular every eigenvalue $\lambda$ of $A$ satisfies $|\lambda|\le\sqrt{pq}<p$.
\end{lemma}
\begin{proof}
Squaring the block matrix of $T_U$ gives diagonal blocks
$q^2+\norm g^2$ and $gg^*+A^2$. The rank-one operator $gg^*$
has trace $\norm g^2$. Consequently, the trace formula gives
\begin{align*}
 q^2+2\norm g^2+\Tr(A^2)
 &=\Tr(T_U^2)
  =\int_{[0,1]^2}U(x,y)^2\dd x\dd y\\
 &\le\int_{[0,1]^2}U(x,y)\dd x\dd y=q,
\end{align*}
where the inequality uses $0\le U\le1$. Subtracting $q^2$ and
using $q-q^2=pq$ proves \eqref{eq:HS}. Since $A$ is self-adjoint,
$\Tr(A^2)$ is the sum of its squared eigenvalues. Thus each
eigenvalue satisfies $\lambda^2\le\Tr(A^2)\le pq$.
Finally $q<p$ gives $\sqrt{pq}<p$.
\end{proof}

\paragraph{Proof strategy after the matrix reduction.}
The value $q$ is the spectral threshold suggested by the desired bound.
\Cref{lem:isolation} will show that the bound holds when every
eigenvalue is at most $q$, and that at most one eigenvalue can exceed
$q$ under the density assumption. Denote that eigenvalue by $\Lambda$.
After isolating it, the proof proceeds as follows.
\begin{enumerate}[label=\arabic*.,leftmargin=*]
 \item The excursion expansion bounds the cycle density using the
 odd moment of $A$ and a quadratic expression in $g$
 (\Cref{sec:excursions}).
 \item Splitting $g$ into its component along the exceptional
 eigenfunction and its orthogonal component reduces this bound to
 two spectral levels, $\Lambda$ and $M<q$
 (\Cref{sec:isolate}).
 \item Kernel positivity forces these two coupling components to
 account for $\Lambda-q$. This is the only step that uses the
 eigenfunction itself (\Cref{sec:compensation}).
 \item A comparison for truncated squares converts that constraint
 into the required bound on every odd moment. Taylor's formula
 explains why the same comparison works for all lengths
 (\Cref{sec:integration,sec:scalar,sec:completion}).
\end{enumerate}
The restriction $p\le2q$ is also used in the coefficient comparison
in \Cref{sec:compensation} and in the lower-cutoff argument in
\Cref{sec:scalar}. Isolating one eigenvalue is therefore the
first use of the density assumption, not its only use.

\paragraph{Relation to the dense proof.}
In the companion paper, adding the cycle expansions for $W$ and $U$
cancels their odd spectral moments. The remaining comparison with a
path-density correction is reduced by Dirichlet averaging to one
spectral variable. Beta and Gamma integral representations then
reduce positivity to an odd-to-even Gamma moment inequality.
Deleting an edge from the complementary cycle gives the one-colour
bound. Here we instead retain the odd spectral moment and compensate
for its exceptional eigenvalue directly. The strengthened
extended-commonality inequality used in the dense proof can fail
below $2/3$, as the companion paper shows by a two-part graphon.
Thus the one-colour bound on the present interval requires a
different comparison.

\section{Excursions: an exact finite counting identity}\label{sec:excursions}
The block decomposition separates the constant and mean-zero
components of a function. In this section, we expand $T_W^m$ by
inserting $I=\Pi_{\mathrm c}+\Pi_0$ between its factors. Each term
specifies which component is retained after each step. We group
these terms by the number of passages from the constants to $E_0$
and back, and prove that the sum over each positive number of such
passages is nonnegative.

A \emph{word} is a sequence of the symbols $\mathrm c,0$, standing
for the projections $\Pi_{\mathrm c},\Pi_0$. For trace expressions
the word is cyclic, so its first and last positions are consecutive.
In a word containing both symbols, an \emph{excursion} is a maximal
nonempty block of $0$'s, with a $\mathrm c$ before and after it.
The number of $0$'s counts the positions in $E_0$, rather than the
number of transitions within $E_0$. Thus a block of $j+1$ zeros
has $j$ internal transitions and two transitions between subspaces.

\begin{lemma}[Projection words and excursion weights]\label{lem:words}
For every integer $m\ge3$,
\begin{equation}\label{eq:words}
 \Tr(T_W^m)=\sum_{\omega\in\{\mathrm c,0\}^m}
 \Tr\bigl(\Pi_{\omega_1}T_W\Pi_{\omega_2}T_W\cdots\Pi_{\omega_m}T_W\bigr).
\end{equation}
The constant-only word contributes $p^m$; the mean-zero-only word
contributes $\Tr((-A)^m)$.

An excursion with $j\ge0$ internal transitions in $E_0$, followed by $k\ge0$ constant-to-constant transitions, has length $j+k+2$ and weight
\[
 p^k\ip{g}{(-A)^j g}.
\]
\end{lemma}
\begin{proof}
Insert $I=\Pi_{\mathrm c}+\Pi_0$ between consecutive copies of $T_W$ in
the trace and expand to obtain \eqref{eq:words}. The pure words
use only the diagonal blocks $p$ and $-A$. An excursion leaves
the constant direction through $-g$, makes $j$ transitions through
$-A$, returns through $-g^*$, and is followed by $k$ factors of
$p$. The two minus signs on departure and return cancel, giving
the stated weight. For example, the segment $\mathrm c\,0\,0\,\mathrm c$
corresponds to
\[
 \Pi_{\mathrm c}T_W\Pi_0T_W\Pi_0T_W\Pi_{\mathrm c}
   =\ip g{(-A)g}\Pi_{\mathrm c}.
\]
The two mean-zero positions give one internal factor of $-A$.
More generally,
\[
 \Pi_{\mathrm c}T_W(\Pi_0T_W)^{j+1}\Pi_{\mathrm c}
  =\ip g{(-A)^jg}\Pi_{\mathrm c},
 \qquad \Pi_{\mathrm c}T_W\Pi_{\mathrm c}=p\Pi_{\mathrm c}.
\]
Following the excursion by $k$ constant-to-constant transitions
therefore multiplies its weight by $p^k$.
\end{proof}
For $\ell\ge2$, define the sum of these weights at total length $\ell$ by
\begin{equation}\label{eq:e}
 e_\ell=\sum_{j=0}^{\ell-2}p^{\ell-2-j}\ip{g}{(-A)^j g}.
\end{equation}

\begin{lemma}[Cyclic excursion expansion]\label{lem:exc}
For every integer $m\ge3$,
\begin{equation}\label{eq:exactexc}
 \Tr(T_W^m)=p^m+\Tr((-A)^m)
 +\sum_{r=1}^{\lfloor m/2\rfloor}\frac mr
 \sum_{\substack{\ell_1+\cdots+\ell_r=m\\\ell_i\ge2}}
 e_{\ell_1}\cdots e_{\ell_r}.
\end{equation}
Every $e_\ell$ is nonnegative. Consequently, for odd $m\ge3$,
\begin{equation}\label{eq:oneexc}
 t(C_m,W)\ge p^m-\Tr(A^m)+m\ip{g}{k_m(A)g},
 \qquad k_m(z)=\frac{p^{m-1}-z^{m-1}}{p+z}.
\end{equation}
The function $k_m(A)$ can equivalently be defined by the polynomial
$\sum_{j=0}^{m-2}p^{m-2-j}(-A)^j$.
\end{lemma}
\begin{proof}
Fix a number $r\ge1$ of excursions. Each excursion uses at least
two transitions, so $r\le\lfloor m/2\rfloor$. For the $i$th
excursion, let $j_i$ be its number of internal mean-zero transitions
and $k_i$ the number of constant-to-constant transitions following
its return. Set $\ell_i=j_i+k_i+2$. The length and weight are
\[
 \ell_1+\cdots+\ell_r=m,\qquad
 \prod_{i=1}^r p^{k_i}\ip g{(-A)^{j_i}g}.
\]
The product formula follows by collapsing each excursion to a
scalar multiple of the rank-one projection $\Pi_{\mathrm c}$ and
using $\Tr(\Pi_{\mathrm c})=1$.

To sum these weights, give the positions of a cyclic word the labels
$1,\ldots,m$ and mark a departure from $\mathrm c$ to $0$.
Each word with $r$ excursions has exactly $r$ possible marks.
Starting at its mark determines an ordered tuple
$(j_1,k_1,\ldots,j_r,k_r)$ of nonnegative integers.
Conversely, choose such a tuple with total length $m$, and choose
one of the $m$ labelled positions for the marked departure. Fill
the cyclic positions with the specified excursions and constant
steps. This constructs one marked word, and the construction is
inverse to extracting the tuple and marked position from a word.

For fixed lengths $\ell_1,\ldots,\ell_r$, summing over the choices
$j_i+k_i=\ell_i-2$ gives
\[
 \sum_{\substack{j_i,k_i\ge0\\j_i+k_i=\ell_i-2\ (1\le i\le r)}}
 \prod_{i=1}^r p^{k_i}\ip g{(-A)^{j_i}g}
 =\prod_{i=1}^r\left(
   \sum_{j_i=0}^{\ell_i-2}p^{\ell_i-2-j_i}\ip g{(-A)^{j_i}g}
   \right)
 =\prod_{i=1}^r e_{\ell_i}.
\]
Summing the weights of the marked words in the two ways therefore gives
\[
 r\,\bigl(\text{total weight of words with $r$ excursions}\bigr)
 =m\sum_{\substack{\ell_1+\cdots+\ell_r=m\\\ell_i\ge2}}
 e_{\ell_1}\cdots e_{\ell_r}.
\]
Dividing by $r$, summing over $r$, and adding the two pure words
proves \eqref{eq:exactexc}. If a word is periodic, different
placements may give the same unmarked sequence of symbols. They
still give distinct marked words, since the marks occupy different
labelled positions. Thus the factor $m/r$ does not assume that
every word has $m$ distinct rotations.

To prove nonnegativity, choose an orthonormal eigenbasis $(\phi_i)$
of $A$, with eigenvalues $\lambda_i$, and put $g_i=\ip g{\phi_i}$.
Then
\[
 \ip g{(-A)^jg}=\sum_i g_i^2(-\lambda_i)^j,
 \qquad
 \sum_{j=0}^{\ell-2}p^{\ell-2-j}(-\lambda_i)^j
  =\frac{p^{\ell-1}-(-\lambda_i)^{\ell-1}}{p+\lambda_i}.
\]
Substituting in the definition \eqref{eq:e} of $e_\ell$ gives
\begin{equation}\label{eq:epos}
 e_\ell=\sum_i g_i^2\frac{p^{\ell-1}-(-\lambda_i)^{\ell-1}}{p+\lambda_i}\ge0.
\end{equation}
By \Cref{lem:HS}, $|\lambda_i|<p$. Thus $p+\lambda_i>0$ and
$p^{\ell-1}-(-\lambda_i)^{\ell-1}>0$, so every $e_\ell$ is
nonnegative. Each product in the expansion \eqref{eq:exactexc}
is therefore nonnegative. For odd $m$, the term with $r=1$ is
$m e_m=m\ip g{k_m(A)g}$, and
$\Tr((-A)^m)=-\Tr(A^m)$. Discarding all terms with $r\ge2$
gives the one-excursion lower bound \eqref{eq:oneexc}.
\end{proof}

\begin{remark}
Individual projection words may have negative weight. Nonnegativity is proved only after the excursion lengths and constant transitions have been summed into the coefficients $e_\ell$. It is these aggregates, not individual words, that are discarded.
\end{remark}

\begin{lemma}[The first two odd excursion expansions]\label{lem:short-expansions}
The cycle densities at lengths three and five satisfy
\begin{align*}
 t(C_3,W)&=p^3-\Tr(A^3)
       +3\bigl(p\norm g^2-\ip g{Ag}\bigr),\\
 t(C_5,W)&=p^5-\Tr(A^5)+5\ip g{k_5(A)g}
       +5\norm g^2\bigl(p\norm g^2-\ip g{Ag}\bigr).
\end{align*}
\end{lemma}
\begin{proof}
In \Cref{lem:exc}, the only excursion count for $m=3$ is
$r=1$. Since $e_3=p\norm g^2-\ip g{Ag}$, this gives the first
identity. For $m=5$, the count $r=1$ contributes
$5\ip g{k_5(A)g}$, while $r=2$ contributes
\[
 \frac52(e_2e_3+e_3e_2)
   =5\norm g^2\bigl(p\norm g^2-\ip g{Ag}\bigr).
\]
In the second identity the last term is the aggregate of the
two-excursion words. It is nonnegative because
$\ip g{Ag}\le\norm A\norm g^2<p\norm g^2$ when $g\ne0$.
\end{proof}
The one-excursion inequality keeps exactly the terms quadratic in $g$;
the expansion proves that the omitted higher-order aggregates cannot
decrease the cycle density.

\section{Isolating the exceptional eigenvalue}\label{sec:isolate}
The one-excursion bound separates the contribution that can lower
the cycle density, namely $\Tr(A^m)$, from a nonnegative expression
in $g$. If every eigenvalue of $A$ is at most $q$, the
Hilbert--Schmidt bound already controls this odd moment. We first
prove this case, and then reduce the remaining case to the
exceptional eigenvalue and a bound on the rest of the spectrum.
\begin{lemma}[Isolation of the exceptional eigenvalue]\label{lem:isolation}
There is at most one eigenvalue of $A$ greater than $q$, counted
with multiplicity. If there is none, then the bound
\eqref{eq:target} holds for every odd $m\ge3$.
\end{lemma}
\begin{proof}
Two eigenvalues greater than $q$ would give
\[
 \Tr(A^2)>2q^2\ge pq,
\]
contrary to \eqref{eq:HS}; the second inequality uses $q\ge1/3$.

If every eigenvalue is at most $q$, then for odd $m$,
\[
 \lambda^m\le q^{m-2}\lambda^2
\]
for every eigenvalue $\lambda$. Indeed, for $0\le\lambda\le q$
this follows from $\lambda^{m-2}\le q^{m-2}$; for $\lambda<0$
the left-hand side is negative and the right-hand side is
nonnegative. Summing over the eigenvalues gives
\[
 \Tr(A^m)\le q^{m-2}\Tr(A^2)\le pq^{m-1}.
\]
The last term of \eqref{eq:oneexc} is $m\ip{g}{k_m(A)g}=me_m\ge0$ by \eqref{eq:epos}, so \eqref{eq:oneexc} gives $t(C_m,W)\ge p^m-pq^{m-1}$.
\end{proof}

It remains to consider a simple eigenvalue $\Lambda>q$. Choose a real unit eigenfunction $\phi$ and write
\begin{equation}\label{eq:split}
 A\phi=\Lambda\phi,\qquad
 \gamma=\ip g\phi,\qquad g=\gamma\phi+g_s,\qquad
 g_s\perp\one,\phi,\qquad
 s=\norm{g_s},\qquad M=\sqrt{pq-\Lambda^2}.
\end{equation}
Thus $\gamma$ measures how the degree fluctuation couples to the
exceptional eigenfunction, whereas $s$ measures its coupling to all
remaining eigenfunctions together. The number $M$ is an upper bound on
their absolute eigenvalues, not necessarily an eigenvalue itself.
Let $A_s$ be the restriction of $A$ to $\{\one,\phi\}^\perp\cap E$.
\begin{lemma}[Bounds on the remaining spectrum]\label{lem:remaining-spectrum}
The remaining operator and the two coupling components satisfy
\begin{equation}\label{eq:safeHS}
 \Tr(A_s^2)+2(\gamma^2+s^2)\le M^2.
\end{equation}
Every eigenvalue of $A_s$ therefore lies in $[-M,M]$, and
\begin{equation}\label{eq:ordering}
 0\le M<q<\Lambda<\frac12<p.
\end{equation}
\end{lemma}
\begin{proof}
Because $A$ is self-adjoint and $A\phi=\Lambda\phi$, the
orthogonal complement of $\phi$ in $E_0$ is invariant under $A$.
Thus $A_s$ is well-defined and self-adjoint. Since
$\Tr(A^2)=\Lambda^2+\Tr(A_s^2)$ and
$\norm g^2=\gamma^2+s^2$, the Hilbert--Schmidt bound gives
\eqref{eq:safeHS}. Each eigenvalue of $A_s$ has square at most
$\Tr(A_s^2)\le M^2$. To compare $M$ with $q$, use $\Lambda>q$
and $p=1-q$ to obtain
\[
 M^2=pq-\Lambda^2<pq-q^2=q(1-2q)\le q^2.
\]
The last inequality is equivalent to $q\ge1/3$. Also,
$\Lambda^2\le pq=p(1-p)<1/4$, since $p>1/2$.
Together with $q<\Lambda$ and $p>1/2$, these inequalities give
the ordering \eqref{eq:ordering}.
\end{proof}

\begin{lemma}[Reduction to two spectral levels]\label{lem:two-levels}
For every odd $m\ge3$, define
\begin{align}\label{eq:E}
 E_m={}&\frac{m\gamma^2}{p+\Lambda}(p^{m-1}-\Lambda^{m-1})
 +\frac{ms^2}{p+M}(p^{m-1}-M^{m-1})\notag\\
 &-\Lambda^m-M^m+pq^{m-1}.
\end{align}
Then
\begin{equation}\label{eq:Ebound}
 t(C_m,W)-(p^m-pq^{m-1})\ge E_m+2M^{m-2}(\gamma^2+s^2).
\end{equation}
\end{lemma}
\begin{proof}
In the orthogonal decomposition
$E=\R\one\oplus\R\phi\oplus(E\cap\{\one,\phi\}^\perp)$,
\begin{equation}\label{eq:threeblocks}
 T_W=\begin{pmatrix}
 p&-\gamma&-g_s^*\\
 -\gamma&-\Lambda&0\\
 -g_s&0&-A_s
 \end{pmatrix}.
\end{equation}
There is no transition between the exceptional direction and the remaining spectral subspace. Consequently a nonzero excursion stays in one of these two subspaces. An exceptional excursion with $b\ge1$ nonconstant positions and $m-b\ge1$ constant positions has, for each rotation, weight
\[
 p^{m-b-1}\gamma^2(-\Lambda)^{b-1}.
\]
There is exactly one departure in such a word, so the marked-word
count gives $m$ placements. Summing over $b$ gives
\[
 m\sum_{b=1}^{m-1}p^{m-b-1}\gamma^2(-\Lambda)^{b-1}
  =m\gamma^2\sum_{j=0}^{m-2}p^{m-2-j}(-\Lambda)^j
  =m k_m(\Lambda)\gamma^2.
\]
The same calculation in the remaining subspace gives
$m\ip{g_s}{k_m(A_s)g_s}$. Thus the two terms separate according
to which invariant subspace the excursion visits.

We next give a lower bound for the remaining-subspace term.
Since $m-1$ is even, for $0\le z<p$,
\[
 k_m(-z)=\frac{p^{m-1}-z^{m-1}}{p-z}
  \ge\frac{p^{m-1}-z^{m-1}}{p+z}=k_m(z).
\]
The numerator is positive and $0<p-z\le p+z$, which justifies
the inequality. With $n=m-1$, the quotient rule also gives
\[
 k_m'(z)=-\frac{p^n+npz^{n-1}+(n-1)z^n}{(p+z)^2}<0.
\]
Thus $k_m(\lambda)\ge k_m(M)$ for $|\lambda|\le M$. Since $A=\Lambda\oplus A_s$ on $\R\phi\oplus(\{\one,\phi\}^\perp\cap E)$, the polynomial $k_m(A)$ splits in the same way, and
\[
 \ip{g}{k_m(A)g}=k_m(\Lambda)\gamma^2+\ip{g_s}{k_m(A_s)g_s}
 \ge k_m(\Lambda)\gamma^2+k_m(M)s^2,
\]
To justify the inequality, let $(\psi_i)$ be an orthonormal
eigenbasis of $A_s$, with eigenvalues $\mu_i\in[-M,M]$, and write
$g_s=\sum_i a_i\psi_i$. Then
\[
 \ip{g_s}{k_m(A_s)g_s}
  =\sum_i a_i^2k_m(\mu_i)
  \ge k_m(M)\sum_i a_i^2=k_m(M)s^2.
\]
Hence
\begin{equation}\label{eq:couplelower}
 m\ip{g}{k_m(A)g}\ge m k_m(\Lambda)\gamma^2+m k_m(M)s^2.
\end{equation}
The words that never visit the constant direction contribute
$-\Lambda^m-\Tr(A_s^m)$. For every eigenvalue $\mu_i$ of $A_s$,
oddness of $m$ and $|\mu_i|\le M$ give
$\mu_i^m\le M^{m-2}\mu_i^2$. Therefore the bound
\eqref{eq:safeHS} on the remaining squared eigenvalues gives
\begin{align}\label{eq:noreturn}
 \Tr(A^m)
 &\le\Lambda^m+M^{m-2}\Tr(A_s^2)\notag\\
 &\le\Lambda^m+M^m-2M^{m-2}(\gamma^2+s^2).
\end{align}
Combining the one-excursion bound \eqref{eq:oneexc}, the coupling
bound \eqref{eq:couplelower}, and the moment bound
\eqref{eq:noreturn} gives \eqref{eq:Ebound}.
\end{proof}
We will prove $E_m\ge0$ without using the extra nonnegative term in \eqref{eq:Ebound}.
The terms in this inequality have a direct interpretation. The
quantities $\Lambda^m+M^m$ bound the odd spectral moment, and
$pq^{m-1}$ is the amount allowed by the theorem. The two remaining
terms are the one-excursion contributions associated with $\gamma$
and $g_s$. Since $M<\Lambda$ and $k_m$ decreases on $[0,p)$, a unit
of squared coupling in the remaining subspace gives at least as large
an excursion contribution as a unit in the exceptional direction.
Keeping these two components separate is useful in the next step.

\begin{lemma}[The one-excursion coefficient as a derivative]\label{lem:coupling-derivative}
For a single spectral level $z\in[0,p)$ and a squared coupling $r\ge0$,
put
\[
 B_z(r)=\begin{pmatrix}p&-\sqrt r\\-\sqrt r&-z\end{pmatrix}.
\]
For every odd $m\ge3$,
\[
 \left.\frac{d}{dr}\Tr(B_z(r)^m)\right|_{r=0}
   =\frac{m}{p+z}\bigl(p^{m-1}-z^{m-1}\bigr)=m k_m(z),
\]
where the derivative at zero is the right derivative.
\end{lemma}
\begin{proof}
The characteristic polynomial is $(\theta-p)(\theta+z)-r$,
so the eigenvalues and their derivatives at zero are
\[
 \theta_\pm(r)=\frac{p-z\pm\sqrt{(p+z)^2+4r}}2,
 \qquad \theta_+'(0)=\frac1{p+z},\quad
 \theta_-'(0)=-\frac1{p+z}.
\]
Since $\theta_+(0)=p$ and $\theta_-(0)=-z$, differentiation gives
\[
 \left.\frac{d}{dr}\Tr(B_z(r)^m)\right|_{r=0}
   =m p^{m-1}\theta_+'(0)+m(-z)^{m-1}\theta_-'(0)
   =m k_m(z).
\]
\end{proof}
Coupling separates the positive and negative eigenvalues. Since
$p>z$, their combined odd moment increases to first order by precisely
the one-excursion coefficient. This calculation explains the
coefficient; the lower bound for the full operator follows from the
exact excursion expansion, not from an approximation in $r$.

\section{Compensation forced by the eigenfunction}\label{sec:compensation}
The preceding section reduced the cycle bound to $E_m\ge0$, but
the Hilbert--Schmidt inequality alone does not give a sufficient
lower bound on its two coupling terms. We now use the constraints
$0\le U\le1$ to relate $\gamma$ and $s$ to $\Lambda-q$.

The comparison is between the exceptional eigenfunction $\phi$ and
its pointwise absolute value. Kernel nonnegativity implies that
passing from $\phi$ to $|\phi|$ cannot decrease the quadratic form
of $T_U$. Unlike $\phi$, the function $|\phi|$ has a nonzero
constant component. The corresponding diagonal entry of $T_U$ is
$q<\Lambda$, so the off-diagonal terms must compensate for this
decrease. Resolving $|\phi|$ into its constant, exceptional, and
remaining components distinguishes the two ways this can happen.

\begin{lemma}[Coupling forced by an exceptional eigenvalue]\label{lem:comp}
One has $M>0$. Write $\gamma_+=\max\{\gamma,0\}$.
The sign of $\phi$ may be chosen so that
\begin{equation}\label{eq:naturalcomp}
 \Lambda-q\le
 \frac{1-2\Lambda}{\sqrt{\Lambda(1-\Lambda)}}\gamma_+
 +\frac{s^2}{\Lambda-M}.
\end{equation}
In particular,
\begin{equation}\label{eq:comp}
 \Lambda-q\le c\gamma_++\frac{s^2}{\Lambda-M},
 \qquad c=\sqrt{\frac{M^2(p-\Lambda)}{\Lambda pq}}.
\end{equation}
\end{lemma}
The first coefficient comes from the positive and negative parts of
the eigenfunction. The slightly larger coefficient $c$ gives a weaker
inequality that is convenient for the later comparison of squares.
The proof separates these two assertions.
\begin{proof}
We first obtain two comparisons involving $|\phi|$, then separate
the two components of $g$, and finally compare the coefficients in
the two stated inequalities. All function decompositions below
are in the finite-dimensional space $E$.

\smallskip
\noindent\emph{The absolute-value comparison.}
Changing $\phi$ to $-\phi$ changes the sign of
$\ip{|\phi|}\phi$. Choose the sign so that this inner product is
nonnegative, and set
\[
 a=\ip\one{|\phi|}=\int_0^1|\phi(x)|\dd x,
 \qquad b=\ip{|\phi|}\phi\ge0.
\]
The orthogonal decomposition of $|\phi|$ is
\[
 |\phi|=a\one+b\phi+k,\qquad k\perp\one,\phi.
\]
Since $\norm\phi=1$, we have $a>0$ and
$a^2+b^2+\norm k^2=\norm{|\phi|}^2=1$.
The positive and negative parts satisfy
$\phi=\phi_+-\phi_-$ and $|\phi|=\phi_++\phi_-$.
Because $\int\phi=0$, their integrals are equal, so
\[
 \int_0^1\phi_+(x)\dd x
  =\int_0^1\phi_-(x)\dd x=\frac a2.
\]
For any nonnegative $f,h$, the kernel bounds imply
\[
 0\le\ip f{T_Uh}
   =\int_{[0,1]^2}f(x)U(x,y)h(y)\dd x\dd y
   \le\left(\int_0^1f\right)\left(\int_0^1h\right).
\]
Using this with $f,h\in\{\phi_+,\phi_-\}$ gives
\begin{align*}
 \Lambda=\ip\phi{T_U\phi}
 &=\ip{\phi_+}{T_U\phi_+}+\ip{\phi_-}{T_U\phi_-}
     -2\ip{\phi_+}{T_U\phi_-}\\
 &\le\left(\int\phi_+\right)^2
      +\left(\int\phi_-\right)^2.
\end{align*}
Consequently,
\begin{equation}\label{eq:alower}
 \Lambda\le\left(\int\phi_+\right)^2+\left(\int\phi_-\right)^2=\frac{a^2}{2}.
\end{equation}
The same expansion, with the cross term added rather than
subtracted, gives
\[
 \ip{|\phi|}{T_U|\phi|}-\ip\phi{T_U\phi}
   =4\ip{\phi_+}{T_U\phi_-}\ge0.
\]
Thus
\begin{equation}\label{eq:abs}
 \ip{|\phi|}{T_U|\phi|}\ge\ip\phi{T_U\phi}=\Lambda.
\end{equation}

\smallskip
\noindent\emph{A constraint on the spectral parameters.}
The largest eigenvalue of $A$ is $\Lambda$, since every other
eigenvalue is at most $M<\Lambda$. Put
$w=|\phi|-a\one\in E_0$, so $\norm w^2=1-a^2$.
Using the block form, the inequality
$\ip w{Aw}\le\Lambda\norm w^2$, and Cauchy--Schwarz in
the absolute-value comparison \eqref{eq:abs}, we obtain
\[
 \Lambda\le qa^2+2a\norm g\sqrt{1-a^2}+\Lambda(1-a^2),
\]
so
\[
 (\Lambda-q)a\le2\norm g\sqrt{1-a^2}.
\]
The left-hand side is positive because $\Lambda>q$ and $a>0$.
Thus $a<1$ and $\norm g>0$. We may square and divide by
$4(1-a^2)>0$. The function $x/(1-x)$ is increasing on $[0,1)$,
and the bound \eqref{eq:alower} gives $a^2\ge2\Lambda$. Hence
\[
 \norm g^2\ge\frac{(\Lambda-q)^2a^2}{4(1-a^2)}
 \ge\frac{\Lambda(\Lambda-q)^2}{2(1-2\Lambda)}.
\]
Together with $2\norm g^2\le M^2$, this first gives $M>0$.
Multiplying by $1-2\Lambda>0$ gives
\begin{align}\label{eq:admiss}
 0&\le M^2(1-2\Lambda)-\Lambda(\Lambda-q)^2\notag\\
  &=(p-\Lambda)\bigl(q(1-\Lambda)-\Lambda^2\bigr).
\end{align}
For the factorisation, substitute $M^2=pq-\Lambda^2$ and use
$p+q=1$; the expanded polynomial on either side is
$\Lambda^3+(q-p)\Lambda^2-q(1+p)\Lambda+pq$.
Since $p>\Lambda$, it follows that $q(1-\Lambda)-\Lambda^2\ge0$.

\smallskip
\noindent\emph{Separating the two coupling components.}
The block decomposition gives
\[
 \ip{|\phi|}{T_U|\phi|}
  =qa^2+2ab\gamma+2a\ip{g_s}{k}
       +\Lambda b^2+\ip k{A_s k}.
\]
Use $\ip k{A_s k}\le M\norm k^2$ and subtract
$\Lambda(a^2+b^2+\norm k^2)=\Lambda$ in the absolute-value
comparison \eqref{eq:abs}. This yields
\[
 (\Lambda-q)a^2+(\Lambda-M)\norm k^2
 \le2ab\gamma+2a\ip{g_s}{k}.
\]
Put $r=\norm k/a$. After division by $a^2$, Cauchy--Schwarz and
completion of the square give
\[
 \Lambda-q\le\frac{2b\gamma}{a}+2sr-(\Lambda-M)r^2
 =\frac{2b\gamma}{a}+\frac{s^2}{\Lambda-M}
   -(\Lambda-M)\left(r-\frac{s}{\Lambda-M}\right)^2.
\]
Therefore
\begin{equation}\label{eq:compintermediate}
 \Lambda-q\le\frac{2b\gamma}{a}+\frac{s^2}{\Lambda-M}.
\end{equation}
This is why the two components enter differently. The component
$b\phi$ of $|\phi|$ has the same eigenvalue as $\phi$, so its
contribution is linear in $\gamma$. The remaining component
satisfies $\ip k{A_s k}\le M\norm k^2$; completing the square
against the gap $\Lambda-M$ produces
the quadratic contribution $s^2/(\Lambda-M)$.

If $\gamma\le0$, then $2b\gamma/a\le0$ and $\gamma_+=0$,
so the coupling bound \eqref{eq:compintermediate} gives both
conclusions of the lemma.

Suppose $\gamma>0$. The pointwise kernel bounds give
\[
 T_U\phi=T_U\phi_+-T_U\phi_-
  \le T_U\phi_+\le\left(\int\phi_+\right)\one
  =\frac a2\one.
\]
The exceptional eigenvalue equation and the block form give
$T_U\phi=\gamma\one+\Lambda\phi$. Also,
\[
 b=\ip{|\phi|}\phi=\norm{\phi_+}^2-\norm{\phi_-}^2,
 \qquad \norm{\phi_+}^2+\norm{\phi_-}^2=1,
\]
so $\norm{\phi_+}^2=(1+b)/2$. Taking the inner product of
$T_U\phi\le(a/2)\one$ with $\phi_+\ge0$ therefore gives
\[
 \frac{a\gamma}{2}+\frac{\Lambda(1+b)}2
  =\ip{\phi_+}{T_U\phi}
  \le\frac a2\int\phi_+=\frac{a^2}{4}.
\]
Multiplying by four and using $a^2+b^2\le1$, we obtain
\[
 2a\gamma+2\Lambda(1+b)\le a^2\le1-b^2.
\]
In particular,
\[
 2\Lambda(1+b)\le1-b^2=(1-b)(1+b),
 \qquad 0\le b\le1-2\Lambda.
\]
Since
$\frac{d}{db}(b^2/(1+b))=b(2+b)/(1+b)^2\ge0$ for $b\ge0$,
the bounds $a^2\ge2\Lambda(1+b)$ and $b\le1-2\Lambda$ give
\[
 \frac{b^2}{a^2}\le\frac{b^2}{2\Lambda(1+b)}
 \le\frac{(1-2\Lambda)^2}{4\Lambda(1-\Lambda)}.
\]
Consequently \eqref{eq:compintermediate} gives the sharper estimate
\begin{equation}\label{eq:sharpercomp}
 \Lambda-q\le\frac{1-2\Lambda}{\sqrt{\Lambda(1-\Lambda)}}\gamma_+
 +\frac{s^2}{\Lambda-M}.
\end{equation}
This proves the natural coupling inequality \eqref{eq:naturalcomp}.

\smallskip
\noindent\emph{Comparison of the two coefficients.}
Write $\varepsilon=\Lambda-q$. The identity
\begin{align}\label{eq:parametercert}
 &(1-\Lambda)M^2(p-\Lambda)-pq(1-2\Lambda)^2\notag\\
 &\quad=\varepsilon\Bigl[(1-2\Lambda)(q-p\Lambda)
 +\varepsilon\bigl(q(1-\Lambda)-\Lambda^2\bigr)\Bigr]
 \ge0
\end{align}
follows by expansion with $p+q=1$ and $M^2=pq-\Lambda^2$. Its sign uses \eqref{eq:admiss} and
\[
 q-p\Lambda\ge q-\frac p2=\frac{3q-1}{2}\ge0.
\]
After division by $\Lambda pq(1-\Lambda)>0$, it states
\[
 \frac{(1-2\Lambda)^2}{\Lambda(1-\Lambda)}\le c^2.
\]
Thus \eqref{eq:sharpercomp} implies \eqref{eq:comp}.
\end{proof}

\begin{lemma}[An exceptional eigenvalue with zero direct coupling]\label{lem:zero-coupling-example}
Partition $[0,1]$ into $B_1,B_2,B_0$ of measures
$\alpha,\alpha,1-2\alpha$, where $\sqrt{1/6}<\alpha<1/2$.
Let $U$ be one on $B_1^2\cup B_2^2$ and zero elsewhere. Put
$W=1-U$ and use the space of functions constant on these three
parts for $E$. Then $A$ has eigenvalues
$\Lambda=\alpha>q$ and $\mu=\alpha(1-2\alpha)$. The coupling
components satisfy
\[
 \gamma=0,\qquad s^2=2\alpha^3(1-2\alpha),\qquad
 \frac{s^2}{\Lambda-\mu}=\Lambda-q.
\]
\end{lemma}
\begin{proof}
The edge density and a unit mean-zero eigenfunction are
\[
 q=2\alpha^2\in(1/3,1/2),\qquad
 \phi=\frac{\one_{B_1}-\one_{B_2}}{\sqrt{2\alpha}},\qquad
 A\phi=\alpha\phi.
\]
Thus $\Lambda=\alpha>q$ is exceptional. The degree function
$d_U(x)=\int_0^1 U(x,y)\dd y$ equals $\alpha$ on $B_1\cup B_2$
and zero on $B_0$. The degree fluctuation is
constant on $B_1\cup B_2$, so its inner product with $\phi$ vanishes:
\[
 \gamma=0,\qquad g=g_s.
\]
Its squared norm is
\begin{align*}
 s^2&=2\alpha(\alpha-q)^2+(1-2\alpha)q^2\\
 &=2\alpha^3(1-2\alpha)^2+4\alpha^4(1-2\alpha)
   =2\alpha^3(1-2\alpha)>0.
\end{align*}
To compute the other eigenvalue, use the values of $g=d_U-q$:
\[
 g=(\alpha-q)\one_{B_1\cup B_2}-q\one_{B_0},
 \qquad T_Ug=\alpha(\alpha-q)\one_{B_1\cup B_2}.
\]
The mean of $T_Ug$ is $2\alpha^2(\alpha-q)$, so subtracting
that mean gives
\begin{align*}
 Ag&=\alpha(1-2\alpha)(\alpha-q)\one_{B_1\cup B_2}
       -2\alpha^2(\alpha-q)\one_{B_0}\\
   &=\alpha(1-2\alpha)g,
\end{align*}
where the second equality uses $q=2\alpha^2$ and
$\alpha-q=\alpha(1-2\alpha)$. Since $g\ne0$ and $g\perp\phi$,
the vectors $g$ and $\phi$ span the two-dimensional space $E_0$.
Its two eigenvalues are therefore $\Lambda=\alpha$ and
$\mu=\alpha(1-2\alpha)$. Finally,
\[
 \frac{s^2}{\Lambda-\mu}
  =\frac{2\alpha^3(1-2\alpha)}{2\alpha^2}
 =\alpha-2\alpha^2=\Lambda-q.
\]
\end{proof}
The remaining-subspace term alone accounts for the exceptional
eigenvalue. Geometrically, $\phi$ distinguishes the two cliques,
whereas $|\phi|$ distinguishes their union from $B_0$. The degree
fluctuation is in this second direction. This explains why a bound
using only $\gamma$ would miss a necessary part of the argument.

\section{Removing the cycle length by positive integration}\label{sec:integration}
The coupling constraint is independent of $m$, whereas the lower
bound $E_m$ contains powers whose degrees grow with the cycle
length. We now express all of these powers through the same family
of functions. This reduces $E_m\ge0$ for every odd $m\ge5$ to
one inequality indexed by a cutoff $t\in[0,p]$.

For a continuously differentiable function $f$ on $[0,p]$, define
$F_f(z)=z^2f(z)$ and the linear expression
\begin{align}\label{eq:moment-functional}
 \mathcal L(f)={}&\frac{\gamma^2}{p+\Lambda}
       \bigl(F_f'(p)-F_f'(\Lambda)\bigr)
 +\frac{s^2}{p+M}\bigl(F_f'(p)-F_f'(M)\bigr)\notag\\
 &+pq f(q)-\Lambda^2 f(\Lambda)-M^2 f(M).
\end{align}
The last three terms compare spectral masses: the two masses
$\Lambda^2$ and $M^2$ sum to $pq$, and the reference expression
places all of this mass at $q$. The first two terms are exactly
the excursion correction for this comparison.

\begin{lemma}[Positive representation by truncated squares]\label{lem:taylor-mixture}
If $f\in C^3([0,p])$ and $f(0)=f'(0)=f''(0)=0$, then
\begin{equation}\label{eq:taylor-mixture}
 f(z)=\frac12\int_0^p f'''(t)(z-t)_+^2\,dt
 \qquad(0\le z\le p).
\end{equation}
If $f'''\ge0$, the weights in this representation are nonnegative.
\end{lemma}
\begin{proof}
The three vanishing initial values allow us to integrate $f'''$
three times without constant terms. Reversing the order of integration
on the region $0\le t\le b\le a\le z$ gives
\begin{align*}
 f(z)&=\int_0^z\int_0^a\int_0^b f'''(t)\dd t\dd b\dd a\\
 &=\int_0^z f'''(t)
       \left(\int_t^z\int_t^a\dd b\dd a\right)\dd t
  =\frac12\int_0^z f'''(t)(z-t)^2\dd t.
\end{align*}
The exchange is valid because $f'''$ is continuous and the region
is compact. The inner double integral is
$\int_t^z(a-t)\dd a=(z-t)^2/2$.
The integrand in \eqref{eq:taylor-mixture} vanishes for $t>z$,
so this is the claimed formula. The last assertion follows from
$f'''(t)\ge0$.
\end{proof}
Thus, whenever $f'''\ge0$, the function $f$ is a positive average of
truncated squares. In particular this applies to $f(z)=z^{m-2}$
for $m\ge5$. This is the reason for testing truncated squares:
their positivity under $\mathcal L$ implies positivity for every
power needed here. No assertion about arbitrary convex functions
is required.

For $t,z\ge0$, define
\begin{equation}\label{eq:D}
 D_t(z)=\frac{d}{dz}\bigl(z^2(z-t)_+^2\bigr)
       \qquad(t,z\ge0).
\end{equation}
Also define
\begin{align}\label{eq:K}
 \mathcal K(t)={}&
 \frac{\gamma^2}{p+\Lambda}\bigl(D_t(p)-D_t(\Lambda)\bigr)
 +\frac{s^2}{p+M}\bigl(D_t(p)-D_t(M)\bigr)\notag\\
 &+pq(q-t)_+^2-\Lambda^2(\Lambda-t)_+^2-M^2(M-t)_+^2.
\end{align}
\begin{lemma}[From cutoff comparisons to powers]\label{lem:power-mixture}
For $t,z\ge0$, one has
$D_t(z)=2z(z-t)_+(2z-t)$ and
$\mathcal L((z-t)_+^2)=\mathcal K(t)$, where $z$ is the
argument of the function supplied to $\mathcal L$.
For every odd integer $m\ge5$, put
$C_m=(m-2)(m-3)(m-4)/2$. Then
\begin{equation}\label{eq:Eint}
 E_m=C_m\int_0^p t^{m-5}\mathcal K(t)\,dt.
\end{equation}
Consequently, if $\mathcal K(t)\ge0$ on $[0,p]$, then $E_m\ge0$
for every odd $m\ge5$.
\end{lemma}
\begin{proof}
For $z\le t$, both $z^2(z-t)_+^2$ and its derivative vanish.
For $z\ge t$, the product rule gives
\[
 D_t(z)=2z(z-t)^2+2z^2(z-t)=2z(z-t)(2z-t).
\]
Substitution in the definition of $\mathcal L$ gives
$\mathcal L((z-t)_+^2)=\mathcal K(t)$.
For $f(z)=z^{m-2}$, we have $F_f'(z)=mz^{m-1}$, so
$\mathcal L(z^{m-2})=E_m$. Substituting
$f'''(t)=2C_m t^{m-5}$ in Taylor's formula
\eqref{eq:taylor-mixture} gives
\begin{equation}\label{eq:powerint}
 z^{m-2}=C_m\int_0^p t^{m-5}(z-t)_+^2\,dt.
\end{equation}
Multiplying by $z^2$ and differentiating gives
\begin{equation}\label{eq:powerintD}
 mz^{m-1}=C_m\int_0^p t^{m-5}D_t(z)\,dt.
\end{equation}
Differentiation under the integral is valid because
$z^2(z-t)_+^2$ and its $z$-derivative are continuous on the compact
square $[0,p]^2$. The four evaluation points $M,q,\Lambda,p$
belong to this interval by the ordering \eqref{eq:ordering}.
Apply the power representation \eqref{eq:powerint} to the last
three terms of $E_m$, and its differentiated form
\eqref{eq:powerintD} to the two excursion terms. Linearity gives
\begin{align*}
 C_m\int_0^p t^{m-5}\mathcal K(t)\dd t
 ={}&\frac{\gamma^2}{p+\Lambda}
       \bigl(mp^{m-1}-m\Lambda^{m-1}\bigr)\\
 &+\frac{s^2}{p+M}\bigl(mp^{m-1}-mM^{m-1}\bigr)\\
 &+pq\,q^{m-2}-\Lambda^2\Lambda^{m-2}-M^2M^{m-2}
 =E_m.
\end{align*}
This is the identity \eqref{eq:Eint}. Since $C_m>0$ and
$t^{m-5}\ge0$, nonnegativity of $\mathcal K$ implies $E_m\ge0$.
\end{proof}
It therefore remains to prove $\mathcal K(t)\ge0$ for every cutoff $t\in[0,p]$. This is a stronger assertion than positivity of the integrals alone, and it contains no cycle-length parameter.

\section{The length-independent scalar lemma}\label{sec:scalar}
We first normalize the spectral levels and the coupling constraint
by the exceptional eigenvalue.
\begin{lemma}[Normalized parameters]\label{lem:normalization}
For the parameters and sign choice in \Cref{lem:comp}, put
\begin{equation}\label{eq:normalizeapplication}
 P=\frac p\Lambda,\qquad x=\frac q\Lambda,\qquad
 y=\frac M\Lambda,\qquad h=1-x,
\end{equation}
and
\begin{equation}\label{eq:normalizecoupling}
 u=\frac{c\gamma_+}{\Lambda},\qquad
 v=\frac{s^2}{\Lambda(\Lambda-M)}.
\end{equation}
Then
\[
 0<y<x<1<P,\qquad Px=1+y^2,\qquad P\le2x,
 \qquad u,v\ge0,\qquad u+v\ge h.
\]
\end{lemma}
\begin{proof}
The ordering \eqref{eq:ordering} and $M>0$ give
$0<y<x<1<P$, and the definitions give $u,v\ge0$.
Dividing the coupling bound \eqref{eq:comp} by $\Lambda$ gives
$u+v\ge h$. Finally,
\[
 Px=\frac{pq}{\Lambda^2}
   =\frac{\Lambda^2+M^2}{\Lambda^2}=1+y^2,
 \qquad P=\frac p\Lambda\le\frac{2q}{\Lambda}=2x.
\]
\end{proof}
The inequality $u+v\ge h$ says that the normalized excess must be accounted for
by the two coupling components. The exceptional component has a cost
quadratic in $u$, because an excursion contains $\gamma^2$; the other
has a cost linear in $v$, because $v$ already contains $s^2$.
The following lemma treats these quantities as arbitrary real
parameters satisfying just the stated relations. Its conclusion is
the cutoff comparison after division by $\Lambda^4$.
\begin{lemma}[Truncated-square inequality]\label{lem:scalar}
Suppose
\begin{equation}\label{eq:normalized}
 0<y<x<1<P,\qquad Px=1+y^2,\qquad P\le2x.
\end{equation}
Set $h=1-x$, and let $u,v\ge0$ satisfy $u+v\ge h$. Put
\[
 A(t)=\frac{D_t(P)-D_t(1)}{P+1},\qquad
 B(t)=\frac{D_t(P)-D_t(y)}{P+y}.
\]
Then, for $0\le t\le P$,
\begin{align}\label{eq:Phi}
 \Phi(t)={}&\frac{1+y^2}{y^2(P-1)}A(t)u^2+(1-y)B(t)v\notag\\
 &+(1+y^2)(x-t)_+^2-(1-t)_+^2-y^2(y-t)_+^2
 \ge0.
\end{align}
\end{lemma}
\begin{proof}
There are four spectral levels, $y<x<1<P$. Above $y$, the term
$y^2(y-t)_+^2$ has vanished, and two coefficient bounds will reduce
the assertion to completions of squares. Below $y$, all terms are
quadratic polynomials in $t$; we will prove that their sum decreases
towards its already-controlled value at $y$. These are the two parts
of the proof.

For fixed $t\ge0$, $D_t(z)$ is nonnegative and nondecreasing
in $z\ge0$. Indeed, it vanishes for $z\le t$, is continuous
at $z=t$, and for $z>t$,
\[
 \frac{d}{dz}D_t(z)=12(z-t)^2+12t(z-t)+2t^2\ge0.
\]
In particular $A(t),B(t)\ge0$.

\smallskip
\noindent\emph{Two coefficient bounds.}
The bounds needed above the cutoff $y$ are
\begin{align}
 A(t)&\ge P-1 &&(0\le t\le1),\label{eq:Abound}\\
 (1-y)B(t)&\ge2(1-t) &&(y\le t\le1).\label{eq:Bbound}
\end{align}
The first makes the exceptional contribution at least
$(1+y^2)u^2/y^2$. The second makes the remaining-subspace
contribution at least $2(1-t)v$. These are precisely the
coefficients that allow $u+v\ge h$ to be used in a square.
For the first bound, $t\le1<P$ implies that the polynomial
formula for $D_t$ applies at both $1$ and $P$. Hence
\[
 A(t)=\frac{4(P^3-1)-6t(P^2-1)+2t^2(P-1)}{P+1}.
\]
Differentiating this polynomial for $0<t<1$ gives
\[
 A'(t)=\frac{2(P-1)}{P+1}\bigl(2t-3(P+1)\bigr)<0.
\]
Thus
\[
 A(t)\ge A(1)=\frac{2P(P-1)(2P-1)}{P+1}\ge P-1;
\]
For the last inequality, subtract $P-1$ from the value at $1$:
\[
 A(1)-(P-1)
  =\frac{(P-1)(4P^2-3P-1)}{P+1}
  =\frac{(P-1)^2(4P+1)}{P+1}\ge0.
\]

For \eqref{eq:Bbound}, set $r=1-t$, so $0\le r\le1-y$. Here $D_t(y)=0$ and
\[
 B(t)=2\frac{P}{P+y}(P-t)(2P-t).
\]
The relation $Px=1+y^2$ and $x<1$ give $P>1+y^2$.
For fixed $y$ and $t\in[y,1]$, the three factors
$P/(P+y)$, $P-t$, and $2P-t$ are nonnegative and nondecreasing
in $P\ge1+y^2$. Therefore
\begin{align*}
 \frac{(1-y)B(t)}2-r
 &\ge\frac{(1-y)(1+y^2)(1+y^2-t)(2+2y^2-t)}{1+y+y^2}-r\\
 &=\frac{(1-y)(1+y^2)(r+y^2)(r+1+2y^2)-r(1+y+y^2)}{1+y+y^2}.
\end{align*}
The denominator is positive. The numerator has the following
nonnegative decomposition:
\begin{align}\label{eq:Bcert}
 &(1-y)(1+y^2)(r+y^2)(r+1+2y^2)-r(1+y+y^2)\notag\\
 &\quad=(1-y)(1+y^2)(r-y)^2\notag\\
 &\qquad\quad+y^2r\bigl((1-y)^2+y^2(2-3y+2y^2)\bigr)\notag\\
 &\qquad\quad+2y^4(1+y^2)(1-y-r)\ge0.
\end{align}
To check the signs, use $0<y<1$, $0\le r\le1-y$, and
$2-3y+2y^2=2(y-3/4)^2+7/8>0$. Each of the three terms on the
right of \eqref{eq:Bcert} is nonnegative. Thus
$(1-y)B(t)/2-r\ge0$, which is the second coefficient bound
\eqref{eq:Bbound}.

\smallskip
\noindent\emph{Cutoffs at or above $y$.}
There are three intervals, according to whether the cutoff has
passed $x$ or $1$.
\begin{enumerate}[label=(\roman*)]
\item Suppose $y\le t\le x$. Write $r=1-t$, so $x-t=r-h$. The reference
mass at $x$ still contributes. Apply the coefficient bounds
\eqref{eq:Abound}--\eqref{eq:Bbound}, and then use $v\ge h-u$:
\begin{align}\label{eq:middlesquares}
 \Phi(t)&\ge\frac{1+y^2}{y^2}u^2+2r(h-u)+(1+y^2)(r-h)^2-r^2\notag\\
 &=\left(y(r-h)-\frac uy\right)^2+(u-h)^2\ge0.
\end{align}
The substitution $v\ge h-u$ is valid even when $u>h$, since its
coefficient $2r$ is nonnegative. Thus there is no need for an
additional case according to which coupling component is larger.
\item Suppose $x\le t\le1$. With $r=1-t$, we have $0\le r\le h$,
and the reference term vanishes. Since $(1+y^2)/y^2\ge1$,
the same bounds give
\begin{align}\label{eq:uppersquare}
 \Phi(t)&\ge u^2+2r(h-u)-r^2\notag\\
 &=(u-r)^2+2r(h-r)\ge0.
\end{align}
\item Suppose $1\le t\le P$. All three truncated-square terms
vanish, as do $D_t(1)$ and $D_t(y)$. Hence
\[
 \Phi(t)=D_t(P)\left(
  \frac{1+y^2}{y^2(P-1)(P+1)}u^2
   +\frac{1-y}{P+y}v\right)\ge0.
\]
\end{enumerate}
These three cases prove $\Phi(t)\ge0$ for every $t\in[y,P]$.

\smallskip
\noindent\emph{A consequence of the density restriction.}
For $t\le y$, the negative term $-y^2(y-t)^2$ reappears, so the
preceding squares alone do not apply. Here we use $P\le2x$ to
bound how far the reference level $x$ can be from $1$.
From $P\le2x$ and $Px=1+y^2$,
\[
 y^2\le2x^2-1=1-4h+2h^2.
\]
In particular $x>1/\sqrt2>1/2$, and
\[
 (2x-1)^2-y^2\ge2(1-x)^2\ge0.
\]
Hence $y\le2x-1$, or $h\le(1-y)/2$. Substituting this in the preceding inequality,
\[
 4h\le1-y^2+2h^2
 \le(1-y)(1+y)+\frac{(1-y)^2}{2}
 =\frac{(1-y)(3+y)}2.
\]
Therefore
\begin{equation}\label{eq:hbound}
 h\le\frac{(1-y)(3+y)}8.
\end{equation}
Set, for the rest of this proof,
\[
 H=(1+y^2)h-y^2(1-y).
\]
For $t\le y$, the derivative of the three non-coupling terms,
divided by two, is
\begin{align*}
 &-(1+y^2)(x-t)+(1-t)+y^2(y-t)\\
 &\qquad=1+y^3-(1+y^2)x
   =(1+y^2)h-y^2(1-y)=H.
\end{align*}
In particular, this derivative is independent of $t$.
If $H\le0$, those terms already
have the derivative sign that we need. If $H>0$, the coupling
terms must supply a derivative of magnitude at least $2H$.
Since $h>0$, \eqref{eq:hbound} implies
\begin{align}\label{eq:Hbound}
 \frac{H}{(1-y)h}
 &=\frac{1+y^2}{1-y}-\frac{y^2}{h}\notag\\
 &\le\frac{1+y^2}{1-y}
       -\frac{8y^2}{(1-y)(3+y)}\notag\\
 &=\frac{3+4y-y^2}{3+y}\notag\\
 &=\frac32-\frac{(1-y)(3-2y)}{2(3+y)}\le\frac32.
\end{align}

\smallskip
\noindent\emph{Monotonicity for cutoffs below $y$.}
On $[0,y]$, all positive parts in the formula \eqref{eq:Phi}
can be removed. Here
\[
 D_t(z)=4z^3-6z^2t+2zt^2\qquad(z\in\{P,1,y\}),
\]
and therefore
\[
 A''(t)=\frac{4(P-1)}{P+1},\qquad
 B''(t)=\frac{4(P-y)}{P+y}.
\]
The resulting polynomial $\Phi$ is convex because
\begin{equation}\label{eq:convex}
 \Phi''(t)=\frac{4(1+y^2)}{y^2(P+1)}u^2
 +\frac{4(1-y)(P-y)}{P+y}v\ge0.
\end{equation}
The second derivatives of the three non-coupling terms cancel:
$2(1+y^2)-2-2y^2=0$. Only the nonnegative coupling terms remain
in \eqref{eq:convex}. For a convex function, the derivative at the
right endpoint is its largest derivative on the interval. It
therefore suffices to prove $\Phi'(y^-)\le0$.
For the endpoint derivative, the same polynomial formulas give
\begin{align*}
 A'(y^-)&=-\frac{2(P-1)}{P+1}\bigl(3(P+1)-2y\bigr),\\
 B'(y^-)&=-\frac{2(P-y)(3P+y)}{P+y}.
\end{align*}
Multiplying these derivatives by their coefficients in $\Phi$
and using the derivative $2H$ of the non-coupling terms yields
\begin{align}\label{eq:derivative}
 \frac12\Phi'(y^-)=H
 &-\frac{1+y^2}{y^2}\left(3-\frac{2y}{P+1}\right)u^2\notag\\
 &-\frac{(P-y)(3P+y)}{P+y}(1-y)v.
\end{align}
We claim that this is nonpositive. If $H\le0$, both coefficients
of the subtracted terms are positive: $P>1>y$ gives
$3-2y/(P+1)>2$ and $(P-y)(3P+y)/(P+y)>0$.
Thus $\Phi'(y^-)/2\le H\le0$ in this case.
Suppose now that $H>0$, so
\begin{equation}\label{eq:Hpositive}
 h>\frac{y^2(1-y)}{1+y^2}.
\end{equation}
First, since $P>1>y$,
\begin{equation}\label{eq:firstcost}
 \frac{1+y^2}{y^2}\left(3-\frac{2y}{P+1}\right)
 \ge\frac{2(1+y^2)}{y^2}\ge\frac{2(1-y)}h.
\end{equation}
For the other coefficient, \eqref{eq:Hpositive} gives
\[
 P=\frac{1+y^2}{1-h}\ge(1+y^2)(1+h)
 \ge1+2y^2-y^3=:P_0.
\]
The function
\[
 \frac{(P-y)(3P+y)}{P+y}=3P-5y+\frac{4y^2}{P+y}
\]
increases in $P\ge y$: its derivative is $3-4y^2/(P+y)^2\ge2$. Here $P_0>y$. At $P=P_0$, the following identity proves that this function is at least two:
\begin{align}\label{eq:derivcert}
 &(P_0-y)(3P_0+y)-2(P_0+y)\notag\\
 &\quad=(1-2y)^2+y^2(1-y)(6y^2-5y+3)
 +3y^4(1-y)^2\ge0,
\end{align}
since $6y^2-5y+3=6(y-5/12)^2+47/24>0$. We have shown
\begin{equation}\label{eq:secondcost}
 \frac{(P-y)(3P+y)}{P+y}\ge2\qquad\text{whenever }H>0.
\end{equation}
By \eqref{eq:firstcost}--\eqref{eq:secondcost} and $v\ge h-u$, the sum of the two subtracted terms in \eqref{eq:derivative} is at least
\begin{align}\label{eq:finalsquare}
 \frac{2(1-y)}h u^2+2(1-y)v
 &\ge\frac{2(1-y)}h u^2+2(1-y)(h-u)\notag\\
 &=2(1-y)h\left[\left(\frac uh-\frac12\right)^2+\frac34\right]\notag\\
 &\ge\frac32(1-y)h\ge H,
\end{align}
where the final inequality is \eqref{eq:Hbound}. Thus $\Phi'(y^-)\le0$.

Convexity on the lower interval gives
$\Phi'(t)\le\Phi'(y^-)\le0$ for $0\le t<y$. Hence $\Phi$
is nonincreasing on $[0,y]$, and
$\Phi(t)\ge\Phi(y)\ge0$ there. Together with the three cases
above $y$, this proves the lemma. This argument uses the left
derivative at $y$ of the polynomial valid on $[0,y]$; it does not
require differentiability across the cutoff.
\end{proof}

\section{Application of the scalar lemma and completion}\label{sec:completion}
We now return from the normalized scalar variables to the graphon
parameters. The homogeneity of the cutoff expression supplies the
factor $\Lambda^4$, and the two coupling coefficients become those
in the truncated-square inequality. This is the last comparison
needed before integrating over the cutoff.
\begin{lemma}[Positivity of the cutoff expression]\label{lem:cutoff-positive}
For the exceptional-eigenvalue parameters of the graphon,
$\mathcal K(t)\ge0$ for every $t\in[0,p]$.
\end{lemma}
\begin{proof}
Choose the sign of $\phi$ as in \Cref{lem:comp} and use
the normalized parameters of \Cref{lem:normalization}.
That lemma verifies all the hypotheses of \Cref{lem:scalar}.
Replacing $\phi$ by $-\phi$ changes $\gamma$ into $-\gamma$ and
leaves $g_s$, $s$, and $E_m$ unchanged, since $E_m$ depends on
$\gamma$ only through $\gamma^2$; so this sign choice does not
affect the lower bound \eqref{eq:Ebound}.

Substitute $p=\Lambda P$, $q=\Lambda x$, $M=\Lambda y$,
and $Px=1+y^2$ in the definition of $c$:
\[
 c^2=\frac{\Lambda^2y^2\,\Lambda(P-1)}
              {\Lambda\cdot\Lambda P\cdot\Lambda x}
     =\frac{y^2(P-1)}{Px}
     =\frac{y^2(P-1)}{1+y^2}.
\]
Since $u=c\gamma_+/\Lambda$ and
$v=s^2/(\Lambda^2(1-y))$, we obtain
\[
 \frac{\gamma^2}{\Lambda^2}\ge
 \frac{\gamma_+^2}{\Lambda^2}
 =\frac{u^2}{c^2}
 =\frac{1+y^2}{y^2(P-1)}u^2,\qquad
 \frac{s^2}{\Lambda^2}=(1-y)v.
\]
The explicit formula for $D_t$ in \Cref{lem:power-mixture} gives
\[
 D_{\Lambda t}(\Lambda z)
  =2\Lambda z\,\Lambda(z-t)_+\,\Lambda(2z-t)
  =\Lambda^3D_t(z).
\]
Also $(\Lambda z-\Lambda t)_+^2=\Lambda^2(z-t)_+^2$.
Using these two scaling identities in the definition \eqref{eq:K}
of $\mathcal K$, and then using the coefficient bounds above, gives
\begin{align*}
 \frac{\mathcal K(\Lambda t)}{\Lambda^4}
 ={}&\frac{\gamma^2}{\Lambda^2}A(t)+\frac{s^2}{\Lambda^2}B(t)\notag\\
 &+(1+y^2)(x-t)_+^2-(1-t)_+^2-y^2(y-t)_+^2
 \ge\Phi(t)\ge0
\end{align*}
for every $t\in[0,P]$. Thus $\mathcal K\ge0$ throughout $[0,p]$.
\end{proof}

\begin{proof}[Proof of \Cref{thm:main}]
\Cref{lem:low-density,lem:triangle} handle
$p\le1/2$ and $m=3$. By \Cref{lem:step-reduction}, it
suffices to consider a step graphon with $1/2<p\le2/3$ and an
odd length $m\ge5$.
If $A$ has no eigenvalue greater than $q$, \Cref{lem:isolation}
proves the bound. Otherwise \Cref{lem:cutoff-positive} and
the integral representation in \Cref{lem:power-mixture} give
$E_m\ge0$. The lower bound in \Cref{lem:two-levels} then yields
\[
 t(C_m,W)\ge p^m-pq^{m-1}.
\]
The same-density step approximation passes this bound to arbitrary
graphons, completing the proof.
\end{proof}

\section*{The comparison beyond odd powers}
The scalar argument also gives a comparison for more general functions.
\begin{lemma}[Comparison for functions with nonnegative third derivative]\label{lem:function-comparison}
For the exceptional-eigenvalue parameters of the proof,
let $\mathcal L$ be the linear expression in
\eqref{eq:moment-functional}. Every $f\in C^3([0,p])$ satisfying
\[
 f(0)=f'(0)=f''(0)=0,\qquad f'''(t)\ge0\quad(0\le t\le p)
\]
also satisfies
\[
 \mathcal L(f)=\frac12\int_0^p f'''(t)\mathcal K(t)\,dt\ge0.
\]
\end{lemma}
\begin{proof}
Apply $\mathcal L$ to the representation in
\Cref{lem:taylor-mixture}. The derivative terms may be passed
under the integral as in the proof of \Cref{lem:power-mixture},
giving the claimed identity. Its integrand is nonnegative by
$f'''\ge0$ and \Cref{lem:cutoff-positive}.
\end{proof}
Thus the
truncated-square lemma compares a class of functions, of which the
odd-cycle powers are particular examples. The graphon input is the
coupling constraint; the coefficient bounds and completions of
squares prove that this constraint suffices for the comparison.

\paragraph{Verification note.}
The accompanying file \nolinkurl{verify_complete_odd_cycle_proof.py} checks the algebraic certificates, the normalization, representative exact excursion identities, and the power-integral formula using rational symbolic arithmetic. These checks supplement the mathematical proof; they are not a formalization in a proof assistant.

\bibliographystyle{abbrv}
\bibliography{references}
\end{document}
